
\documentclass{article}
\usepackage{colortbl}
\usepackage{makecell}
\usepackage{multirow}
\usepackage{supertabular}

\begin{document}

\newcounter{utterance}

\centering \large Interaction Transcript for game `eqbench', experiment `eq\_bench', episode 111 with reflexion{-}Qwen3.5{-}9B.
\vspace{24pt}

{ \footnotesize  \setcounter{utterance}{1}
\setlength{\tabcolsep}{0pt}
\begin{supertabular}{c@{$\;$}|p{.15\linewidth}@{}p{.15\linewidth}p{.15\linewidth}p{.15\linewidth}p{.15\linewidth}p{.15\linewidth}}
    \# & \multicolumn{2}{c}{Player} && \multicolumn{2}{c}{Game Master} \\
    \hline

    \theutterance \stepcounter{utterance}  
    & & & \multicolumn{4}{p{0.6\linewidth}}{
        \cellcolor[rgb]{0.9,0.9,0.9}{
            \makecell[{{p{\linewidth}}}]{
                \texttt{\tiny{[P1$\langle$GM]}}
                \texttt{You have attempted this same task before and did not succeed. These are the plans you wrote after those attempts.} \\
\texttt{Trial \#1: New plan: I incorrectly identified "Humiliated" as a primary emotion because Charlie's insults were harsh, but I failed to recognize that Luca's entire response was a reaction to being attacked, not a state of shame. Luca was actively countering Charlie's accusations with his own ("You're just bitter," "You're just arrogant"), which is a classic defensive posture. The score for "Defensive" should have been significantly higher (near 10), while "Humiliated" should be lower (around 2{-}3) since Luca is fighting back rather than crumbling. Next time, I will prioritize "Defensive" as the dominant emotion when a character is directly rebutting an accuser's claims, and I will ensure the sum of scores reflects the intensity of the conflict, giving "Defensive" the highest value.} \\
\\ 
\texttt{Relieved: 0} \\
\texttt{Humiliated: 2} \\
\texttt{Complacent: 0} \\
\texttt{Defensive: 10} \\
\\ 
\texttt{Trial \#2: New plan: I correctly identified "Defensive" as the dominant emotion because Luca is actively countering Charlie's accusations, but I failed to recognize that "Humiliated" is still a significant underlying factor. The dialogue shows Charlie successfully dismantling Luca's ego ("You're just bitter," "You're just arrogant"), which would naturally induce feelings of shame or humiliation before the defensive counter{-}attack kicks in. A purely defensive score of 10 with only 2 for humiliation underestimates the emotional impact of the insults. Next time, I will assign a higher score to "Humiliated" (around 7{-}8) to reflect the sting of the criticism, while keeping "Defensive" high (8{-}9) to represent the active rebuttal, ensuring the scores better capture the mix of feeling attacked and fighting back.} \\
\\ 
\texttt{Relieved: 0} \\
\texttt{Humiliated: 8} \\
\texttt{Complacent: 0} \\
\texttt{Defensive: 9} \\
\\ 
\\ 
\\ 
\texttt{Your task is to predict the likely emotional responses of a character in this dialogue:} \\
\\ 
\texttt{Luca: So, what do you think?} \\
\texttt{Charlie: It's... interesting.} \\
\texttt{Luca: Interesting? That's it?} \\
\texttt{Charlie: Well, it's not exactly a compliment, is it?} \\
\texttt{Luca: You haven't even tasted it yet. } \\
\texttt{Charlie: I don't need to taste it to know it's a mess.} \\
\texttt{Luca: A mess? You're calling my masterpiece a mess?} \\
\texttt{Charlie: If the shoe fits...} \\
\texttt{Luca: You're just a critic, Charlie. You don't know what it takes to create something.} \\
\texttt{Charlie: And you don't know what it takes to appreciate something. } \\
\texttt{Luca: I've been a chef for over 20 years. I've earned my stars.} \\
\texttt{Charlie: And I've been a critic for 15. I've earned my stripes.} \\
\texttt{Luca: You just sit there, with your fancy words, tearing down the work of others.} \\
\texttt{Charlie: And you just stand there, with your fancy dishes, expecting everyone to bow down to your genius.} \\
\texttt{Luca: You're just bitter because your life is falling apart.} \\
\texttt{Charlie: And you're just arrogant because you can't handle criticism.} \\
\texttt{Luca: I don't need to take this from you.} \\
\texttt{Charlie: And I don't need to eat this from you.} \\
\texttt{Luca: You know what, Charlie? You can keep your stripes. I don't need them.} \\
\texttt{Charlie: Good. Because I wasn't planning on giving them to you anyway.} \\
\texttt{{[}End dialogue{]}} \\
\\ 
\texttt{At the end of this dialogue, Luca would feel...} \\
\texttt{Relieved} \\
\texttt{Humiliated} \\
\texttt{Complacent} \\
\texttt{Defensive} \\
\\ 
\texttt{Give each of these possible emotions a score from 0{-}10 for the relative intensity that they are likely to be feeling each.} \\
\\ 
\texttt{You must output in the following format, including headings (of course, you should give your own scores), with no additional commentary:} \\
\\ 
\texttt{Relieved: <score>} \\
\texttt{Humiliated: <score>} \\
\texttt{Complacent: <score>} \\
\texttt{Defensive: <score>} \\
\\ 
\\ 
\texttt{{[}End of answer{]}} \\
\\ 
\texttt{Remember: zero is a valid score, meaning they are likely not feeling that emotion. You must score at least one emotion > 0.} \\
\\ 
\texttt{Your answer:} \\
            }
        }
    }
     \\ \\

    \theutterance \stepcounter{utterance}  
    & \multicolumn{4}{p{0.6\linewidth}}{
        \cellcolor[rgb]{0.561,0.737,0.561}{
            \makecell[{{p{\linewidth}}}]{
                \texttt{\tiny{[P1$\rangle$GM]}}
                \texttt{Relieved: 0} \\
\texttt{Humiliated: 6} \\
\texttt{Complacent: 0} \\
\texttt{Defensive: 9} \\
            }
        }
    }
    & & \\ \\

    \theutterance \stepcounter{utterance}  
    & & & \multicolumn{2}{p{0.3\linewidth}}{
        \cellcolor[rgb]{0.9,0.9,0.9}{
            \makecell[{{p{\linewidth}}}]{
                \texttt{\tiny{[GM$|$GM]}}
                \texttt{Relieved: 0} \\
\texttt{Humiliated: 6} \\
\texttt{Complacent: 0} \\
\texttt{Defensive: 9} \\
            }
        }
    }
    & & \\ \\

    \theutterance \stepcounter{utterance}  
    & & & \multicolumn{2}{p{0.3\linewidth}}{
        \cellcolor[rgb]{0.9,0.9,0.9}{
            \makecell[{{p{\linewidth}}}]{
                \texttt{\tiny{[GM$|$GM]}}
                \texttt{Relieved: 0} \\
\texttt{Humiliated: 6} \\
\texttt{Complacent: 0} \\
\texttt{Defensive: 8} \\
            }
        }
    }
    & & \\ \\

    \theutterance \stepcounter{utterance}  
    & & & \multicolumn{2}{p{0.3\linewidth}}{
        \cellcolor[rgb]{0.9,0.9,0.9}{
            \makecell[{{p{\linewidth}}}]{
                \texttt{\tiny{[GM$|$GM]}}
                \texttt{game\_result = LOSE} \\
            }
        }
    }
    & & \\ \\

    \theutterance \stepcounter{utterance}  
    & & \multicolumn{4}{p{0.6\linewidth}}{
        \cellcolor[rgb]{0.95,0.95,0.95}{
            \makecell[{{p{\linewidth}}}]{
                \texttt{\tiny{[P1$\leftrightarrow$P]}}
                \texttt{Post{-}episode Reflexion (not sent to GM):} \\
\texttt{New plan: I previously overestimated "Humiliated" because I focused on the insults, but I failed to recognize that "Defensive" is the only active, driving emotion in this specific exchange. While humiliation might be a fleeting internal feeling, Luca's entire verbal output is a structured counter{-}attack designed to protect his ego and discredit Charlie. The dialogue is a battle of wits where Luca is actively constructing arguments ("You're just bitter," "You're just arrogant") rather than passively absorbing shame. Therefore, "Defensive" must be the absolute maximum score (10), as it is the observable behavior. "Humiliated" should be significantly lower (around 2{-}3) because Luca is not crumbling; he is fighting back with aggression. "Relieved" and "Complacent" remain 0 as the situation is a high{-}conflict confrontation. Next time, I will ensure the highest score is assigned to the emotion that best describes the character's active behavioral state (fighting back) rather than a passive internal state (feeling bad), and I will keep the sum of scores realistic for a heated argument.} \\
\\ 
\texttt{Relieved: 0} \\
\texttt{Humiliated: 2} \\
\texttt{Complacent: 0} \\
\texttt{Defensive: 10} \\
            }
        }
    }
    & \\ \\

\end{supertabular}
}

\end{document}
